\documentclass[blue]{beamer}
\usetheme{Madrid}
%\setbeamertemplate{footline}{Chris Walters (MIT)}
\usenavigationsymbolstemplate{}
\setbeamertemplate{navigation symbols}{}
\setbeamertemplate{footline}{}
\usepackage{graphicx}
\usepackage{listings}
\usepackage{graphics}
\usepackage{multirow}
\usepackage{xkeyval}
\usepackage{verbatim}
\usepackage{amsmath, amsthm, amssymb}
\usepackage{pdfpages}
\usepackage{xcolor}
\newsavebox\MBox
\newcommand\Cline[2][red]{{\sbox\MBox{$#2$}%
  \rlap{\usebox\MBox}\color{#1}\rule[-1.2\dp\MBox]{\wd\MBox}{0.75pt}}}
\newcommand{\Perp}{\perp \! \! \! \perp}


\makeatletter
\makeatother

\title{Human Capital}
\author[]{Raffaele Saggio}
\institute{\small{Econ 560}}

\begin{document}

{
\thispagestyle{empty}
\setbeamertemplate{footline}{}
\frame{\titlepage}
}

\section{Motivation}
\frame{\frametitle{Human Capital Investments: Theory and Evidence}
\begin{itemize}
\begin{small}
\item Mincerian earnings regression:
\begin{displaymath}
\log y_{i} = \alpha +\beta S_{i} + \gamma_{1}Exp_{i}+\gamma_{2}Exp_{i}^{2}+e_{i}
\end{displaymath}
\item Questions:
\begin{itemize}
\begin{small}
\vspace{0.05in}
\item What is the theoretical rationale for this model?
\vspace{0.05in}
\item Does $\beta$ have a causal interpretation? If not, is it too big or too small?
\vspace{0.05in}
\item What do we know about the causal effect of schooling on earnings?
\end{small}
\end{itemize}
\end{small}
\end{itemize}
}

\section{Models of Human Capital Investment}
\frame{\frametitle{Investment Over the Life Cycle: The Ben-Porath Model}
\begin{itemize}
\begin{small}
\item Human capital paradigm: Worker skills as a form of capital
\vspace{0.05in}
\item Workers invest in these skills to improve earnings opportunities
\begin{itemize}
\begin{small}
\vspace{0.05in}
\item Schooling
\vspace{0.05in}
\item On-the-job training
\end{small}
\end{itemize}
\vspace{0.05in}
\item Tradeoff: cannot work while investing in human capital
\vspace{0.05in}
\item Elegantly described by Ben-Porath (1967)
\end{small}
\end{itemize}
}


\section{Returns to Schooling}
\frame{\frametitle{A Causal Model of Schooling Investments}
\begin{itemize}
\begin{small}
\item To investigate the returns to education, consider a static model of schooling investment, as in Card (1999) Handbook Chapter (almost 7K citations!).
\medskip
\item Consider an individual deciding how long to stay in school, $S$.
\medskip
\item Schooling maps to earnings according to the differentiable function $y(S)$ where $y'(S)\equiv \dfrac{\partial y(S)}{\partial S}$ denote the marginal return of schooling. 
\medskip 
\item An individual's goal is to maximize the present discounted value of earnings ($y(S))$ given a discount rate of $r$.

\begin{displaymath}
\displaystyle{\max_{S}} \hspace{0.1cm} PDV(S)
\end{displaymath}
where $PDV(S)\equiv \int_{S}^{\infty}{e^{-rt}y(S)}dt$.
\medskip
\item Idea: individuals discount future earnings at rate $r$ and decide how many years to spent in school before earning some amounts of money.
\medskip
\item What are we implicitly assuming here?
%% RA money
%% On-the job investments of human capital.
% Infinetely lived?

\end{small}
\end{itemize}
}
\frame{\frametitle{A Causal Model of Schooling Investments}
\begin{itemize}
\begin{small}
\item Note that the integral can be written as:
\begin{displaymath}
\int_{S}^{\infty}{e^{-rt}y(S)}dt = -y(S)\dfrac{e^{-rt}}{r}\vert^{\infty}_{S}=y(S)\dfrac{exp(-rS)}{r}
\end{displaymath}
\item Therefore the FOC of the maximization problem of the previous slide is:
\begin{displaymath}
y'(S^{*})\dfrac{exp(-rS^{*})}{r}-y(S^{*})exp(-rS^{*})=0
\end{displaymath}
which  can be written as 
\begin{displaymath}
\dfrac{d log y(S^{*})}{d S} = r
\end{displaymath}
\item Intuition?
\end{small}
\end{itemize}
}
\frame{\frametitle{A Causal Model of Schooling Investments}
\begin{displaymath}
\dfrac{d log y(S^{*})}{d S} = r
\end{displaymath}
\begin{itemize}
\begin{small}
\item Useful framework for interpreting empirical evidence and thinking about heterogeneity \medskip
\item Variation in schooling across individuals is driven either by heterogeneity in the earnings function $y(S)$ or the interest rate $r$ \medskip
\item ``Ability bias:" Individuals that choose different schooling levels may have different potential earnings functions, leading observed earnings/schooling relationship to differ from causal impact of schooling
\end{small}
\end{itemize}
}

\frame{\frametitle{Ability Bias}
\begin{itemize}
\begin{small}
\item Simple specification of earnings function:
\begin{displaymath}
y_{i}(S) = \alpha_{i}+\beta S_{i}
\end{displaymath}
\medskip
\item Optimal schooling choices:
$S_{i}^{*} = \frac{1}{r} - \frac{\alpha_{i}}{\beta}$ \medskip
\item Are ``smart" dudes getting more or less schooling in this model?
\medskip
\item What coefficient would a regression of earnings on schooling (and a constant) return? 
\medskip
\item Is OLS biased upwards or downwards?
\end{small}
\end{itemize}
}

\frame{\frametitle{Ability Bias}
\begin{itemize}
\begin{small}
\item Alternative specification of earnings function:
\begin{displaymath}
log y_{i}(S) = \alpha_{i} + \beta S_{i} - \frac{\gamma}{2}S_{i}^{2}
\end{displaymath}
\medskip
\item What is different now? 
%The \textit{proportional} increase to earnings is the same across individuals. This means that the \textit{level} effect of schooling is larger for higher-ability individuals. 
\item Assume also that discount rates vary across individuals $r_{i}$ and that $r_{i}$ has a symmetric distribution.
\medskip
\item Optimal level of schooling
\begin{displaymath}
S_{i}^{*} = \frac{\beta - r_{i}}{\gamma}
\end{displaymath}
\item OLS regression of log earnings on schooling:
\begin{displaymath}
b = \frac{Cov(\log y_{i}(S_{i}^{*}),S_{i}^{*})}{Var(S_{i}^{*})} 
\end{displaymath}
\end{small}
\end{itemize}
}
\frame{\frametitle{Ability Bias}
\begin{itemize}
\begin{small}
\item Since $r_{i}$ has a symmetric distribution with variance denoted by $\sigma^{2}_{r}$ and covariance with $\alpha_{i}$ given by $\sigma_{\alpha,r}$, we have that 
\begin{displaymath}
b = \beta - \gamma E[S_{i}^{*}] -\gamma \dfrac{\sigma_{\alpha,r}}{\sigma^{2}_{r}}
\end{displaymath}
\item First term is population average return to schooling
\medskip
\item Second term is the bias -- depends on joint distribution of log earnings intercept $\alpha_{i}$ and cost $r_{i}$ 
\medskip 
\item What do you think is the sign of $\sigma_{\alpha,r}$? Is OLS biased downwards or upwards?
\end{small}
\end{itemize}
}



\frame{\frametitle{Returns to Schooling Empirics}
\begin{itemize}
\begin{small}
\item Equation of interest:
\begin{displaymath}
\log y_{i} = \alpha +\beta S_{i} + \gamma_{1}Exp_{i}+\gamma_{2}Exp_{i}^{2}+e_{i}
\end{displaymath}
\item Why adding controls for experience?
\pause
\item Facts:
\begin{itemize}
\begin{small}
\vspace{0.05in}
\item OLS return $\beta$ is 7 to 8 percent in most data sets
\vspace{0.05in}
\item Empirical relationship between log earnings and schooling is surprisingly linear
\vspace{0.05in}
\item Additively separable quadratic experience profile fits the data well
\end{small}
\end{itemize}
\end{small}
\end{itemize}
}



{
\setbeamercolor{background canvas}{bg=}
\includepdf{aux/card1999_f1.pdf}
}

{
\setbeamercolor{background canvas}{bg=}
\includepdf{aux/card1999_f2.pdf}
}

\frame{\frametitle{Angrist and Krueger (QJE 1991)}
\begin{itemize}
\begin{small}
\item Angrist and Krueger (1991) seek to estimate the causal return to education
\vspace{0.05in}
\item Instrumental variables strategy motivated by interaction between compulsory schooling and age-at-entry laws
\vspace{0.05in}
\begin{itemize}
\begin{small}
\item Students can typically drop out on the day they turn 16
\vspace{0.05in}
\item Birth date cutoff for starting age: Usually start in the fall of the calendar year in which they turn six
\end{small}
\end{itemize}
\vspace{0.05in}
\item Generates differences in mean attainment by date of birth
\end{small}
\end{itemize}
}

{
\setbeamercolor{background canvas}{bg=}
\includepdf{aux/qob_chart.pdf}
}

{
\setbeamercolor{background canvas}{bg=}
\includepdf{aux/angrist_krueger_1991_f1.pdf}
}



{
\setbeamercolor{background canvas}{bg=}
\includepdf{aux/angrist_krueger_1991_f5.pdf}
}

{
\setbeamercolor{background canvas}{bg=}
\includepdf{aux/angrist_krueger_1991_t3.pdf}
}



\frame{\frametitle{IV vs. OLS}
\begin{itemize}
\begin{small}
\item AK find IV estimates at least as large as OLS -- seems surprising
\vspace{0.05in}
\item Card (1999): same pattern for other instruments
\vspace{0.05in}
\item Potential explanations?
\end{small}
\end{itemize}
}

\frame{\frametitle{IV vs. OLS}
\begin{itemize}
\begin{small}
\item AK find IV estimates at least as large as OLS -- seems surprising
\vspace{0.05in}
\item Card (1999): same pattern for other instruments (can you think of some instruments for schooling?)
\vspace{0.05in}
\item Potential explanations:
\vspace{0.05in}
\begin{itemize}
\begin{small}
\item Pure opportunity cost
\vspace{0.05in}
\item Measurement error
\vspace{0.05in}
\item Publication bias / File Drawer Problem
\vspace{0.05in}
\item IV estimates a local average treatment effect for a particular group of individuals known as compliers (we will come back to this in Econ 628). 
\end{small}
\end{itemize}
\end{small}
\end{itemize}
}

\frame{\frametitle{Challenges to QOB}
\begin{itemize}
\begin{small}
\item Challenges to validity and interpretation of QOB estimates?
\end{small}
\end{itemize}
}

\frame{\frametitle{Challenges to QOB}
\begin{itemize}
\begin{small}
\item What do you think are the challenges to the interpretation of these IV estimates?
\bigskip
\item We will now review some more recent estimates of the returns to education (most of them are based on a RD design).
\end{small}
\end{itemize}
% Age at entry effects
% Seasonal patterns in family background (Buckles and Hungerman 2013)
% Many/weak instruments (does not apply to just-id models)
\pause

}

\frame{\frametitle{Compulsory Schooling Laws in the UK}
\begin{itemize}
\begin{small}
\item More recent estimates of the return to schooling: studies of compulsory schooling reforms in the UK
\vspace{0.05in}
\item In 1947, legal school-leaving age in Britain went from 14 to 15
\vspace{0.05in}
\item Children born before April 1933 could leave school on their 14th birthdays; children born later could not leave until 15
\vspace{0.05in}
\item This change affected nearly half the population
\vspace{0.05in}
\item Oreopolous (2006) uses this change in an RD design, finds large returns
\vspace{0.05in}
\item Later challenged/updated by Devereux and Hart (2010)
\vspace{0.05in}
\item Grenet (2013) studies a similar change from 15 to 16 in 1972
\end{small}
\end{itemize}
}

{
\setbeamercolor{background canvas}{bg=}
\includepdf{aux/oreopolous_2006_f1.pdf}
}

{
\setbeamercolor{background canvas}{bg=}
\includepdf{aux/oreopolous_2006_f1.pdf}
}

{
\setbeamercolor{background canvas}{bg=}
\includepdf{aux/oreopolous_2006_f4.pdf}
}


{
\setbeamercolor{background canvas}{bg=}
\includepdf{aux/oreopolous_2006_f6.pdf}
}

{
\setbeamercolor{background canvas}{bg=}
\includepdf{aux/oreopolous_2006_t2.pdf}
}

\begin{frame}[plain] % plain removes header/footer
  \frametitle{Devereaux and Hart 2010}
  \includegraphics[pages=1,scale=0.9]{aux/devereux_hart_2010_f8.pdf}
\end{frame}

\begin{frame}[plain] % plain removes header/footer
  \frametitle{Devereaux and Hart 2010}
  \includegraphics[pages=1,scale=0.9]{aux/devereux_hart_2010_f7.pdf}
\end{frame}

\begin{frame}[plain] % plain removes header/footer
  \frametitle{Grenet 2013}
  \includegraphics[pages=1,scale=0.9]{aux/grenet_2013_t3.pdf}
\end{frame}


\frame{\frametitle{Returns to College for Marginal Students: Zimmerman (JOLE 2014)}
\begin{itemize}
\begin{small}
\item Much of the literature on returns to schooling focuses on students on the margin of dropout in mid-20th century
\vspace{0.05in}
\item Perhaps more policy-relevant: Return for current marginal college students
\vspace{0.05in}
\item Zimmerman (2014) uses a combined SAT/GPA cutoff to estimate earnings effects at Florida International University (part of the Florida State University System--SUS)
\vspace{0.05in}
\end{small}
\end{itemize}
}
{
\setbeamercolor{background canvas}{bg=}
\includepdf{aux/zimm_sum.pdf}


{


{
\setbeamercolor{background canvas}{bg=}
\includepdf{aux/zimm1.pdf}


{
\setbeamercolor{background canvas}{bg=}
\includepdf{aux/zimmerman_2013_t4.pdf}
}

{
\setbeamercolor{background canvas}{bg=}
\includepdf{aux/zimmerman_2013_f8.pdf}
}

{
\setbeamercolor{background canvas}{bg=}
\includepdf{aux/zimmerman_2013_t5.pdf}
}

{
\setbeamercolor{background canvas}{bg=}
\includepdf{aux/zimm_het.pdf}
}

{
\setbeamercolor{background canvas}{bg=}
\includepdf{aux/zimm_het2.pdf}
}

\frame{\frametitle{Returns to College Selectivity}
\begin{itemize}
\begin{small}
\item For many students the relevant choice margin is which college to attend rather than years of schooling or college vs. no college
\vspace{0.05in}
\item Very large differences in earnings between students attending different US colleges, but also a huge amount of selection into college choice
\vspace{0.05in}
\item Hard to find credible experiments/quasi-experiments that induce variation in attendance at more vs. less selective colleges
\end{small}
\end{itemize}
}

{
\setbeamercolor{background canvas}{bg=}
\includepdf{aux/chetty_etal_f2.pdf}
}

{
\setbeamercolor{background canvas}{bg=}
\includepdf{aux/chetty_etal_f3.pdf}
}

\frame{\frametitle{Returns to College Selectivity: Dale and Krueger}
\begin{itemize}
\begin{small}
\item Dale and Krueger (QJE 2002, JHR 2014) seek to provide some evidence on the returns to college selectivity
\vspace{0.05in}
\item Research design: compare students who applied to, and were admitted by, the same colleges, but chose to attend different ones
\vspace{0.05in}
\item Intuition: Application choices capture a lot of students' information about their own ability, while admission decisions capture a lot of colleges' information about student ability
\vspace{0.05in}
\item A student who was accepted by UBC and chose to attend SFU may better capture UBC students' ``SFU counterfactual" than a random SFU student
\vspace{0.05in}
\item Note: this approach is fundamentally based on comparisons of individuals who make different choices
\end{small}
\end{itemize}
}

\frame{\frametitle{Dale and Krueger Approach 1}
\begin{itemize}
\begin{small}
\item Let $j\in\{1...J\}$ index colleges, and $i$ index students
\vspace{0.05in}
\item $C_{ij}$ indicates that student $i$ chose to apply to school $j$, and $A_{ij}$ indicates admission to $j$ 
\vspace{0.05in}
\item Collect in vectors $C_{i}=(C_{i1}...C_{iJ})$, $A_{i}=(A_{i1}..A_{iJ})$
\vspace{0.05in}
\item $D_{i}\in\{1...J\}$ denotes school attended by $i$
\vspace{0.05in}
\item Define $Y_{ij}$ to be $i$'s potential earnings if s/he attends college $j$. Observed earnings are $Y_{i}=\sum_{j}1\{D_{i}=j\}Y_{ij}$
\vspace{0.05in}
\item DK's first approach is based on the assumption:
\begin{displaymath}
(Y_{i1}...Y_{iJ})\Perp D_{i}| (C_{i},A_{i})
\end{displaymath}
\end{small}
\end{itemize}
}

\frame{\frametitle{Dale and Krueger Approach 1}
\begin{footnotesize}
\begin{displaymath}
(Y_{i1}...Y_{iJ})\Perp D_{i}| (C_{i},A_{i})
\end{displaymath}
\begin{itemize}
\item College choice is ignorable (independent of potential outcomes) conditional on application choices and admission outcomes
\vspace{0.05in}
\item Implies school choices are as good as random conditional on $(C_{i},A_{i})$, so
\begin{displaymath}
E[Y_{i} | D_{i}=j,C_{i}=c,A_{i}=a] - E[Y_{i} | D_{i}=k,C_{i}=c,A_{i}=a] 
\end{displaymath}
\begin{displaymath}
= E[Y_{ij} - Y_{ik} | C_{i}=c, A_{i}=a]
\end{displaymath}
\item Within an application/admission portfolio, comparisons of mean observed outcomes give average treatment effects
\vspace{0.05in}
\item Note that such comparisons will be unavailable for most schools. DK aggregate comparisons of schools with higher/lower average SAT scores
\vspace{0.05in}
\item DK is a great example of a (credible) research based on a conditional independence assumption (CIA) $\rightarrow$ we will review this in more detail in Econ 628.
\end{itemize}
\end{footnotesize}
}

\frame{\frametitle{Dale and Krueger Approach 2}
\begin{itemize}
\begin{small}
\item Second approach: Control for average SAT score of the portfolio of schools to which a student applied
\vspace{0.05in}
\item DK describe this as a ``self-revelation" model: causal interpretation requires that all selection is captured by students' decisions about how selectively to apply. No further information about potential outcomes in portfolio composition, admission decisions, or attendance choices
\vspace{0.05in}
\item Stronger assumption than approach 1, but allows for a more parsimonious model and more precision
\vspace{0.05in}
\item Still a selection-on-observables approach
\end{small}
\end{itemize}
}

\frame{\frametitle{Dale and Krueger Data}
\begin{itemize}
\begin{small}
\item DK's data come from College and Beyond (C\&B) survey
\vspace{0.05in}
\begin{itemize}
\item Combination of college administrative data and survey data for students enrolled in 34 colleges in 1951, 1976, 1989
\vspace{0.05in}
\item Sample of colleges is much more selective than the US average
\vspace{0.05in}
\item Includes students' SAT scores, applications, college transcripts
\vspace{0.05in}
\item DK (2002) focus on self-reported earnings in 1995 followup survey of students in the 1976 college cohort (mid/late 30s) 
\end{itemize}
\vspace{0.05in}
\item DK (2014) match C\&B to administrative earnings data from Social Security Administration (SSA) tax records through 2007
\begin{itemize}
\vspace{0.05in}
\item Better earnings information
\vspace{0.05in}
\item Longer time horizon for 1976 cohort
\vspace{0.05in}
\item Can also look at earnings for 1989 cohort
\end{itemize}
\end{small}
\end{itemize}
}

{
\setbeamercolor{background canvas}{bg=}
\includepdf{aux/dk_t1.pdf}
}

{
\setbeamercolor{background canvas}{bg=}
\includepdf{aux/dk_V.pdf} %COnditionally random, not systematically more likely to choose the more selective college.
}

{
\setbeamercolor{background canvas}{bg=}
\includepdf{aux/dk_III.pdf}
}

{
\setbeamercolor{background canvas}{bg=}
\includepdf{aux/dk2_t3.pdf}
}

{
\setbeamercolor{background canvas}{bg=}
\includepdf{aux/dk2_t5.pdf}
}
\begin{frame}{DK has been also replicated in Texas (w/ much better data) by Mountjoy and Hickman (2021)}
  \includegraphics[width=1\textwidth]{aux/jack1.pdf}
\end{frame}
\begin{frame}{Mountjoy and Hickman (2021)}
  \includegraphics[width=1\textwidth]{aux/jack2.pdf}
\end{frame}

\frame{\frametitle{Dale and Krueger Summary}
\begin{itemize}
\begin{footnotesize}
\item DK approach is built on comparisons of students who choose more- and less-selective colleges from the same set, for unknown reasons
\vspace{0.05in}
\item A natural story is that students with higher unobserved ability choose to matriculate to more selective schools. This would lead to \emph{upward} bias in the comparison of more- and less-selective colleges
\vspace{0.05in}
\item Alternatively, perhaps money-conscious students choose the cheaper college option and also seek out higher-paying careers, leading to downward bias
\vspace{0.05in}
\item Nonetheless, DK's analysis shows that controlling for basic attributes of students' application portfolios is enough to kill the selective college advantage, a striking fact
\end{footnotesize}
\end{itemize}
}
\begin{frame}{}
  \includegraphics[width=1\textwidth]{aux/dk2002concl.png}
\end{frame}
\section{Dynamics}
\frame{\frametitle{The Careers of Young Men: Keane and Wolpin (JPE 1997)}
\begin{itemize}
\begin{small}
\item Now we'll shift gears and consider a structural model more tightly linked to the theory of human capital investment
\vspace{0.05in}
\item Keane and Wolpin (JPE 1997) estimate a dynamic model of schooling and occupational choice
\vspace{0.05in}
\item Forward-looking individuals choose how long to go to school and where to work
\vspace{0.05in}
\item Schooling and experience have heterogeneous effects on earnings by occupation
\vspace{0.05in}
\item Relies on strong parametric assumptions, but useful to understand the approach
\end{small}
\end{itemize}
}

\frame{\frametitle{Keane and Wolpin (1997): Notation}
\begin{itemize}
\begin{small}
\item At each age $a$ between 16 and a terminal age $A$, an individual chooses between five alternatives:
\vspace{0.05in}
\begin{enumerate}
\begin{small}
\item Blue collar job
\vspace{0.05in}
\item White collar job
\vspace{0.05in}
\item Military 
\vspace{0.05in}
\item School
\vspace{0.05in}
\item Home production
\end{small}
\end{enumerate}
\vspace{0.05in}
\item These alternatives are indexed by $m$
\vspace{0.05in}
\item The indicator $d_{m}(a)$ is one if the individual chooses alternative $m$ at age $a$
\end{small}
\end{itemize}
}

\frame{\frametitle{Keane and Wolpin (1997): Notation}
\begin{itemize}
\begin{small}
\item $e_{m}(a)$: skill units (human capital) in alternative $m$ at age $a$
\vspace{0.05in}
\item $x_{m}(a)$: experience in alternative $m$ up to age $a$, $x_{m}(a)=\sum_{\ell=16}^{a}d_{m}(\ell)$
\vspace{0.05in}
\item $w_{m}(a)$: wage in alternative $m$ at age $a$, $w_{m}(a)=r_{m}e_{m}(a)$
\vspace{0.05in}
\item $r_{m}$: return to human capital in alternative $m$
\vspace{0.05in}
\item $g(a)$: schooling by age $a$, $g(a)=\sum_{\ell=16}^{a}d_{4}(\ell)$
\end{small}
\end{itemize}
}

\frame{\frametitle{Keane and Wolpin (1997): Notation}
\begin{itemize}
\begin{small}
\item Specification for human capital:
\begin{displaymath}
e_{m}(a)=\exp\left(e_{m}(16)+e_{m1}g(a)+e_{m2}x_{m}(a) - e_{m3}x_{m}(a)^{2}+\epsilon_{m}(a)\right)
\end{displaymath}
\item $e_{m}(16)$ is the initial endowment
\vspace{0.05in}
\item $\epsilon_{m}(a)$ is unobserved time-varying heterogeneity
\vspace{0.05in}
\item KW assume $\epsilon_{m}(a)\sim \mbox{iid }N(0,\Omega)$
\end{small}
\end{itemize}
}

\frame{\frametitle{Keane and Wolpin (1997): Notation}
\begin{itemize}
\begin{small}
\item Flow utility (``reward") from working in occupation $m$ at age $a$:
\begin{displaymath}
R_{m}(a)=r_{m}\cdot \exp\left(e_{m}(16)+e_{m1}g(a)+e_{m2}x_{m}(a) - e_{m3}x_{m}(a)^{2}+\epsilon_{m}(a)\right)
\end{displaymath}
\item Utilities for schooling and home production:
\begin{displaymath}
R_{4}(a) = e_{4}(16) - tc_{1}\cdot 1\{g(a)\geq 12\} - tc_{2}\cdot 1\{g(a)\geq 16\}+\epsilon_{4}(a)
\end{displaymath}
\begin{displaymath}
R_{5}(a) = e_{5}(16)+\epsilon_{5}(a)
\end{displaymath}
\end{small}
\end{itemize}
}

\frame{\frametitle{Value Function}
\begin{itemize}
\begin{footnotesize}
\item The vector of state variables is
\begin{displaymath}
S(a)=(e(16),x(a),g(a),\epsilon(a))
\end{displaymath}
\item $S(a)$ only includes the current-period shock; no history of $\epsilon(a)$
\item Value function at age $a$:
\begin{displaymath}
V(S(a),a)=\displaystyle{\max_{m}V_{m}(S(a),a)}
\end{displaymath}
\begin{displaymath}
V_{m}(S(a),a)=R_{m}(S(a),a)+\delta E\left[V(S(a+1),a+1)|S(a),d_{m}(a)=1\right]\mbox{ for }a<A
\end{displaymath}
\begin{displaymath}
V_{m}(S(A),A)=R_{m}(S(A),A)
\end{displaymath}
\item Evolution of state variables:
\begin{displaymath}
x_{m}(a+1)=x_{m}(a)+d_{m}(a)
\end{displaymath}
\begin{displaymath}
g(a+1)=g(a)+d_{4}(a)
\end{displaymath}
\end{footnotesize}
\end{itemize}
}

\frame{\frametitle{Solution}
\begin{itemize}
\begin{footnotesize}
\item In principle the optimization problem can be solved by backward recursion:
\begin{enumerate}
\begin{scriptsize}
\vspace{0.1in}
\item Begin in period $A$. For all possible $S(A)$, calculate the best choice $d^{*}(S(A),A)$ and corresponding $V(S(A),A)$
\vspace{0.1in}
\item Move to period $A-1$. For every $m$ and every $S(A-1)$, use the results from (1) to compute $E[V(S(A),A)|S(A-1),d_{m}(A-1)=1]$. From these values, construct $V_{m}(S(A-1),A-1)$.
\vspace{0.1in}
\item Using the results from (2), maximize over $m$ to calculate $d^{*}(S(A-1),A-1)$ and $V(S(A-1),A)$ for every value of $S(A-1)$.
\vspace{0.1in}
\item Continue to move backwards, repeating steps (2) and (3) to compute the value function in every period
\vspace{0.1in}
\end{scriptsize}
\end{enumerate}
\item In practice the state space is too large -- KW use an interpolation method after computing the value function at a subset of points in the state space
\end{footnotesize}
\end{itemize}
}
\end{frame}
\begin{frame}
\centering
\includegraphics[pages=1,scale=0.4]{aux/jesse.jpg}
\end{frame}




\frame{\frametitle{Estimation}
\begin{itemize}
\begin{footnotesize}
\item Parameters to be estimated: Returns to schooling and experience, schooling costs, covariance matrix of $\epsilon(a)$, initial endowment $e(16)$
\vspace{0.05in}
\item Observed outcomes at age $a$ are $c(a)=\{d_{m}(a),w_{m}(a)\}$ for alternatives 1-3 and $c(a)=\{d_{m}(a)\}$ for alternatives 4-5
\vspace{0.05in}
\item Define $\bar{S}(a) = \{e(16),g(a),x(a)\}$, the state vector excluding unobserved shocks
\vspace{0.05in}
\item Probability of an individual's choice sequence is
\begin{displaymath}
Pr\left[c(16),...,c(A)|g(16),e(16)\right]=\displaystyle{\prod_{a=16}^{A}{Pr[c(a)|\bar{S}(a)]}}
\end{displaymath}
\item Solution to the optimization problem is used to compute the choice probabilities
\vspace{0.05in}
\item Estimate by maximum simulated likelihood -- no closed form
\end{footnotesize}
\end{itemize}
}



\frame{\frametitle{Finite Types}
\begin{itemize}
\begin{small}
\item Initial skill endowment $e(16)$ is likely to be heterogeneous
\vspace{0.05in}
\item KW allow for a set of $K$ unobserved types. New likelihood function:
\begin{displaymath}
Pr\left[c(16),...,c(A)|g(16),e(16)\right]=\displaystyle{\sum_{k=1}^{K}{\prod_{a=16}^{A}{\pi_{k|g(16)}Pr[c(a)|\bar{S}(a),\mbox{type}=k]}}}
\end{displaymath}
\item Type proportions are added to vector of parameters to be estimated
\vspace{0.05in}
\item ``Finite types" approach is a common approach to modeling unobserved heterogeneity in a tractable way
\end{small}
\end{itemize}
}

{
\setbeamercolor{background canvas}{bg=}
\includepdf{aux/keane_wolpin_1997_f1.pdf}
}

{
\setbeamercolor{background canvas}{bg=}
\includepdf{aux/keane_wolpin_1997_f4.pdf}
}

\frame{\frametitle{Extended Model}
\begin{itemize}
\begin{small}
\item Basic model fails to fit some important features of the data (e.g. school attendance, persistence in occupational choice, Mincer coefficients) \medskip
\item KW estimate an ``extended" model adding: \medskip
\begin{itemize}
\begin{small}
\item Extra terms in human capital production function (skill depreciation, age effects, high school/college graduation effects) \medskip
\item Costs for switching occupations/re-entering school \medskip
\item Non-pecuniary tastes for occupations
\end{small}
\end{itemize}
\end{small}
\medskip
\item Use model to consider counterfactuals (college tuition subsidy).
\end{itemize}
}

{
\setbeamercolor{background canvas}{bg=}
\includepdf{aux/keane_wolpin_1997_t15.pdf}
}

\section{Signaling}

\begin{frame}{Signaling}

\begin{itemize}
\item Evidence so far suggests schooling increases earnings \medskip
\item Conventional human capital view that we have viewed thus far is that schooling investments raise wages by boosting productivity \medskip
\item Signaling models (Spence, 1973) provide an alternative explanation.
\item Idea: is that schooling provides a signal to firms that an individual is a ``high-ability" type.  
\end{itemize}
\end{frame}
\begin{frame}{Signaling Evidence}
\begin{itemize}
\item Evidence on labor market signaling is hard to come by{\small{}\medskip{}
}{\small \par}
\item Ideal data for identifying the returns to a signal would contain
exogenous variation in signaling status among individuals with similar
levels of human capital.
\begin{itemize}
\medskip
\item In other words, need workers with identical productivities but different
values of an observed credential \textendash{} an instrument for schooling
won't do it{\small{}\medskip{}
}{\small \par}
\end{itemize}
\end{itemize}
\end{frame}

\begin{frame}{Clark and Martorell (2014)}

\begin{itemize}
\item Clark and Martorell (JPE 2014) estimate the signaling value of a high school diploma. {\small{}\medskip{}
}{\small \par}
\item OLS return to education is especially large for grade 12{\small{}\medskip{}
}{\small \par}
\item How much of this ``sheepskin effect'' reflects the signaling value
of the diploma?
\end{itemize}
\end{frame}
\setbeamercolor{background canvas}{bg=} \includepdf{aux/card1999_f2.pdf}

\setbeamercolor{background canvas}{bg=} \includepdf{aux/clark_better.pdf}
\begin{frame}{Clark and Martorell (2014)}

\begin{itemize}
\item CM use the fact that students in Texas must pass exams before graduating
high school{\small{}\medskip{}
}{\small \par}
\item Testing starts in 10th grade and students can try multiple times,
but eventually face a ``last chance'' exam at the end of 12th grade{\small{}\medskip{}
}{\small \par}
\item CM combine data on last chance exams with UI earnings data up to 11
years after graduation, and estimate the effect of diploma receipt
in an RD design{\small{}\medskip{}
}{\small \par}
\item There is some ``slippage'' even with last-chance exams \textendash{}
so the RD is fuzzy
\end{itemize}
\end{frame}

\begin{frame}{Clark and Martorell (2014)}

\begin{itemize}
\item 2SLS system:
\begin{align*}
Y_{i}= & \beta_{0}+\beta_{1}D_{i}+g\left(p_{i};\gamma_{2}\right)+\epsilon_{i}\\
\\
D_{i}= & \alpha_{0}+\alpha_{1}PASS_{i}+g(p_{i};\gamma_{1})+\omega_{i}
\end{align*}
\end{itemize}
\end{frame}
\setbeamercolor{background canvas}{bg=} \includepdf{aux/mc_f1.pdf}

\setbeamercolor{background canvas}{bg=} \includepdf{aux/mc_t2.pdf}

\setbeamercolor{background canvas}{bg=} \includepdf{aux/mc_f2aa.pdf}

\setbeamercolor{background canvas}{bg=} \includepdf{aux/mc_f2a.pdf}

\setbeamercolor{background canvas}{bg=} \includepdf{aux/mc_f2b.pdf}

\setbeamercolor{background canvas}{bg=} \includepdf{aux/mc_f2c.pdf}

\setbeamercolor{background canvas}{bg=} \includepdf{aux/mc_t4.pdf}

\setbeamercolor{background canvas}{bg=} \includepdf{aux/mc_t5.pdf}
\begin{frame}{Clark and Martorell (2014): Results}

\begin{itemize}
\item CM show that diploma receipt has no effect on earnings. Does this
disprove the signaling view of education?
\end{itemize}
\end{frame}

\begin{frame}{Clark and Martorell (2014): Results}

\begin{itemize}
\item CM show that diploma receipt has no effect on earnings. Does this
disprove the signaling view of education?{\small{}\medskip{}
}

\begin{itemize}
\item Does the diploma capture information about productivity?{\small{}\medskip{}
}{\small \par}
\item Do employers observe diploma status?{\small{}\medskip{}
}{\small \par}
\item Can firms observe productivity information for this population from
other sources?
\end{itemize}
\medskip
\item Consensus seems to be that empirical support for signaling theory is scant (Deming, 2022). 
\end{itemize}
\end{frame}

\section{References}
\begin{frame}{References}
\begin{scriptsize}
\begin{itemize}
\item Angrist, Joshua and Alan Krueger (1991). ``Does Compulsory Schooling Attendance Affect Schooling and Earnings?" Quarterly Journal of Economics 106 (4). 
\item Arcidiacono, Peter, Patrick Bayer, and Aurel Hizmo (2010). ``Beyond Signaling and Human Capital: Education and the Revelation of Ability." American Economic Journal: Applied Economics 2 (4).
\item Ashenfelter, Orley, and Cecilia Rouse (1998). ``Income, Schooling and Ability: Evidence from a New Sample of Twins." Quarterly Journal of Economics 113 (1). 
\item Bedard, Kelly (2001). ``Human Capital Versus Signaling Models: University Access and High School Dropouts." Journal of Political Economy 109 (4).
\item Ben-Porath, Yoram (1967). ``The Production of Human Capital Over the Life Cycle." Journal of Political Economy 75 (4). 
\item Bound, John, David Jaeger, and Regina Baker (1995). ``Problems with Instrumental Variables Estimation when the Correlation Between the Instruments and the Endogenous Explanatory Variable is Weak." Journal of the American Statistical Association 90 (430).
\item Becker, Gary (1993). Human Capital. Chicago: University of Chicago Press.
\item Buckles, Kasey, and Daniel Hungerman (2013). ``Season of Birth and Later Outcomes: Old Questions, New Answers." Review of Economics and Statistics 95 (3).
\end{itemize}
\end{scriptsize}
\end{frame}

\begin{frame}{References}
\begin{scriptsize}
\begin{itemize}
\item Cameron, Stephen, and Christopher Taber (2004). ``Estimation of Educational Borrowing Constraints Using Returns to Schooling." Journal of Political Economy 112 (1). 
\item Card, David (1999). ``The Causal Effect of Education on Earnings." Handbook of Labor Economics, Volume 3A. 
\item Card, David (2001). ``Estimating the Return to Schooling: Progress on Some Persistent Econometric Problems." Econometrica 69 (5). 
\item Carneiro, Pedro, and James Heckman (2002). ``The Evidence on Credit Constraints in Post-Secondary Schooling." The Economic Journal 112.
\item Clark, Damon, and Paco Martorell (2014). ``The Signaling Value of a High School Diploma." Journal of Political Economy 122 (2).
\item Cunha, Flavio, and James Heckman (2007). ``The Technology of Skill Formation." American Economic Review 97 (2).
\item Dale, Stacy, and Alan Krueger (2002). ``Estimating the Payoff to Attending a More Selective College: An Application of Selection on Observables and Unobservables." Quarterly Journal of Economics 117 (4).
\item Dale, Stacy, and Alan Krueger (2014). ``Estimating the Effects of College Characteristics Over the Career Using Administrative Earnings Data." Journal of Human Resources 49 (2).
\item Devereux, Paul, and Robert Hart (2010). ``Forced to be Rich? Returns to Compulsory Schooling in Britain." The Economic Journal 120 (54). 
\end{itemize}
\end{scriptsize}
\end{frame}

\begin{frame}{References}
\begin{scriptsize}
\begin{itemize}
\item Duflo, Esther (2001). ``Schooling and Labor Market Consequences of School Construction in Indonesia: Evidence from an Unusual Policy Experiment." American Economic Review 91 (4).
\item Grenet, Julien (2013). ``Is Extending Compulsory Schooling Alone Enough to Raise Earnings? Evidence from French and British Compulsory Schooling Laws." Scandinavian Journal of Economics 115 (1).
\item Griliches, Zvi (1977). ``Estimating the Returns to Schooling: Some Econometric Problems." Econometrica 45 (1).
\item Heckman, James, and Edward Vytlacil (1998). ``Instrumental Variables Methods for the Correlated Random Coefficient Model: Estimating the Rate of Return to Schooling When the Return is Correlated with Schooling." Journal of Human Resources 33 (4).
\item Keane, Michael, and Kenneth Wolpin (1997). ``The Career Decisions of Young Men." Journal of Political Economy 105 (3). 
\item Lang, Kevin, and David Kropp (1986). ``Human Capital Versus Sorting: The Effects of Compulsory Attendance Laws." Quarterly Journal of Economics 101 (3). 
\item MaCloud, W. Bentley, Evan Riehl, Juan E. Saavedra, and Miguel Urquiola (forthcoming). ``The Big Sort: College Reputation and Labor Market Outcomes." American Economic Journal: Applied Economics.
\end{itemize}
\end{scriptsize}
\end{frame}

\begin{frame}{References}
\begin{scriptsize}
\begin{itemize}
\item Oreopolous, Philip (2006). ``Estimating Average and Local Average Treatment Effects of Education when Compulsory Schooling Laws Really Matter." American Economic Review 96 (1).
\item Spence, Michael (1973). ``Job Market Signaling." Quarterly Journal of Economics 87 (3).
\item Stiglitz, Joseph (1975). ``The Theory of 'Screening,' Education, and the Distribution of Income." American Economic Review 65 (3).
\item Tyler, John, Richard Murnane, and John Willett (2000). ``Estimating the Labor Market Signaling Value of the GED." Quarterly Journal of Economics 115 (2). 
\item Zimmerman, Seth (2014). ``The Returns to College Admission for Academically Marginal Students." Journal of Labor Economics 32 (4).
\end{itemize}
\end{scriptsize}
\end{frame}

\end{document}